\section{Introduction}
\label{sec:intro}

On-policy self-distillation (\opsd{}) has emerged as an exciting approach for self-improvement in language models \mbox{\citep{zhao2026selfdistilledreasoneronpolicyselfdistillation,shenfeld2026self,hubotter2026reinforcement}}.
In this setting, a single model plays the role of both a student and  teacher.
The teacher is provided additional privileged information, such as a gold solution, a final answer, or environmental feedback.

Thinking models are natural candidates for self-improvement. Post-trained to deliberate at test time \mbox{\citep{openai2024learning,deepseek2025r1,yang2025qwen3}},
they can branch into cases, verify intermediate steps, hedge, backtrack, and recover from errors before committing to an answer \mbox{\citep{arora2025training,gandhi2025cognitive,venhoff2025understanding}}.
Yet existing methods have mostly been studied outside this regime, using instruction-tuned models, short generation budgets, or trained on rollouts with thinking disabled  ~\citep{zhao2026selfdistilledreasoneronpolicyselfdistillation, shenfeld2026self, hubotter2026reinforcement}. This leaves open whether privileged-context self-distillation remains beneficial when the supervised trajectory is itself the long deliberation trace that thinking models rely on at test time.

In this paper, we report a negative result: existing privileged-context on-policy self-distillation methods can degrade thinking models, particularly at long rollout budgets. This degradation is not explained by the short training completion budget alone: under the same short-budget on-policy training setup, vanilla OPD with a larger teacher improves the thinking student, whereas providing the teacher with privileged context reverses the gain. Our analysis suggests a plausible explanation, namely that privileged context may suppress the deliberative behaviors that thinking models rely on at test time.

We establish this failure mode through five linked observations. First, \opsd{} helps instruction-tuned models more reliably than thinking models. Second, under full-solution privileged context, \opsd{} degrades five OpenThoughts-trained thinking models across Qwen3 and OLMo. Third, the harm scales with how much privileged context the teacher receives: final-answer-only context causes milder damage, while full-solution context causes substantially more. Fourth, the gold-demo-conditioned teacher itself does not show the same degradation: it produces much shorter rollouts while preserving high pass@k, whereas the trained student inherits shorter responses without the same longer-budget benefits. Fifth, the degradation is strongest at long rollout budgets and coincides with reduced fork rates, weaker self-correction signal, and fewer deliberation markers in trained students.

Such findings called for mechanistic exploration at the token level. The explanation appears to be a crucial sign reversal in the per-token distillation signal at high-entropy decision points, suggesting that access to privileged context suppresses exploration --e.g.,  branching, reconsideration, and uncertainty-marking moves. Thinking-model rollouts contain many such positions, which we call {\em forking positions} in line with recent work~\citep{bigelow2024forking,lin2024critical,vassoyan2025ignore,zhang2026embarrassingly}. They are often marked lexically by tokens such as \emph{wait}, \emph{hmm}, \emph{but}, and \emph{maybe}. Under vanilla on-policy distillation, these tokens carry positive advantage. Once privileged context is added, their advantage flips negative, and the trained student produces fewer of them at evaluation, even after length normalization.

We refer to this phenomenon as fork suppression: privileged-context self-distillation undermines the very deliberative behaviors that made thinking models natural candidates for self-improvement in the first place.

Section~\ref{sec:method} describes the experimental setup; Section~\ref{sec:results} establishes the empirical pattern; Section~\ref{sec:analysis} analyzes the per-token signal at forking positions; Section~\ref{sec:relatedwork} situates the result among recent self-distillation methods; Section~\ref{sec:Discussion} discusses implications for self-improvement of thinking models.
